% Pursuing Stacks (cotes 134-2 to 134-8): archive cote 134-5, site page 74 (PDF page 75), typescript page 261.
% typescript. Draft compared: ps4.tex:3733-3810.
% AI-transcribed by Claude Opus 5.5 (claude-opus-5-5, adaptive thinking, effort high), Message Batches API, batch ps2 (msgbatch_01NVwr1Yuo8qJf5rykeGBXPK), 2026-11-12; not yet checked by eye.
\scanpage{134-5/74}
\mspage{261}

(Cf pages 213-214  for the conditions Loc 1)  to  Loc 3)\unsure{\hand{.}})

Among all basic localizers satisfying this extra condition,  there
is one coarsest of all,  which we'll call $\underline{W}_{o}$, defined by the condition
\[ (11) \qquad f \hand{\in} \underline{W}_{o} \hand{\Leftrightarrow} \underline{\hand{\pi}}_{o}(f) \text{ is bijective} \]
for any map f in (Cat).  It is clear too that among all basic localizers
there is a finest,  which we'll call $\underline{W}_{\infty}$ ,  and which can be describes as
\[ (12) \qquad \underline{W}_{\infty} = \hand{\bigcap} \underline{W} \hand{,} \quad \text{intersection of the set of all basic localizers in (Cat) .} \]
We'll see in part \unsure{\hand{V}} of the notes that $\underline{W}_{\infty}$ is none else than just the usu-%
al notion of weak equivalence\deleted{,}  we started with\hand{,} at the very beginning of
our reflections  (cf. section 17).  Thus functoriality of (1) with respect
to $\underline{W}$ implies that (1) \xout{is} for arbitrary $\underline{W}$ can be described in terms of the
particular case $\underline{W}_{\infty}$ ,  as the composition
\[ (13) \qquad \mathrm{Cont}(M) \hand{\hookrightarrow} \underline{W}_{\infty}\text{-Asph}(M) \hand{\rightarrow} \underline{W}\text{-Asph}(M) \quad . \]
On the other hand,  the strongest version of injectivity of (1),  for
different $\underline{W}$'s ,  is obtained for $\underline{W}_{o}$ ,  i.e.  taking the map
\[ (14) \qquad \mathrm{Cont}(M) \hand{\hookrightarrow} \underline{W}_{o}\text{-Asph}(M) \ . \]

This last map seems to me of special significance,  because the
two sets it relates correspond to ``\uline{absolute}'' notions  (not depending on
the choice of some $\underline{W}$ ),  and which moreover are \unsure{both} ``\uline{elementary}''\hand{,} in the
sense that they do not depend on anything like consideration of non tri-%
vial homotopy or (co)homology invariants of objects of (Cat)\unsure{\hand{.}}  \over{a}{A}s a matte
% NOTE: the end of this line ("matte...") is cut off at the right edge of the scan
of fact,  the notion of a contractibility structure corresponds to the
algebraic translation of one of the most elementary and intuitive topolo-%
gical notions,  \xout{whereas the} namely contractibility;  whereas the notion
of a $\underline{W}_{o}$-Asphericity structure can be expressed just as ``elementarily'' in
terms of the functor \xout{:}
\[ (15) \qquad \underline{\hand{\pi}}_{o} : (\mathrm{Cat}) \hand{\longrightarrow} (\mathrm{Sets}) \quad , \]
which we may call the ``basic'' connected components-functor,  which is
nothing but the algebraic counterpart of the basic intuitive notion of
connected components of a space.  We'll denote by
\[ (16) \qquad M_{o} = M_{\underline{W}_{o}} \hand{\subset} \mathrm{Ob}\, M \ , \quad \xout{Wx}\ W_{o} \hand{\subset} \mathrm{Fl}(M) \]
the set of $\underline{W}_{o}$-aspheric objects and the set of $\underline{W}_{o}$-weak equivalences,  asso
% NOTE: the end of this line ("asso...") is cut off at the right edge of the scan
ciated to a given contractibility structure $M_{c}$ on M.  The objects of $M_{o}$

\margin{74}
